\section{Part D. Advanced DP and Reconstruction}
\begin{frame}{길이와 실제 LCS는 다르다}
\begin{columns}[T]\column{.48\textwidth}\begin{block}{Length}$dp[m][n]=4$\end{block}
\column{.48\textwidth}\begin{block}{Solution}backtracking으로 \texttt{BCBA}\end{block}\end{columns}
\mission{어떤 문자를 선택했는지 표의 dependency를 거꾸로 추적한다.}
\end{frame}
\begin{frame}{Backtracking 규칙}
\begin{enumerate}\item 문자가 같으면 그 문자를 선택하고 diagonal 이동\item 다르면 큰 DP 값의 방향으로 이동\item 동점이면 위쪽을 우선\end{enumerate}
\begin{alertblock}{Convention}동점 규칙에 따라 다른 올바른 LCS가 나올 수 있다.\end{alertblock}
\end{frame}
\begin{frame}{LCS Backtracking Trace}
\centering
\begin{tikzpicture}[
  lcs table/.style={matrix of nodes,row sep=1.2pt,column sep=1.2pt,
    nodes={minimum width=8.4mm,minimum height=5.2mm,font=\normalsize,anchor=center},
    ampersand replacement=\&},
  lcs header/.style={font=\normalsize\bfseries,text=AlgoBlue},
  trace move/.style={-Latex,very thick,draw=AlgoBlueTwo},
  chosen/.style={draw=AlgoOrange,very thick,rounded corners=1pt,inner sep=1pt}
]
\matrix[lcs table] (lcs) {
  |[lcs header]|{} \& |[lcs header]|{$\boldsymbol{\epsilon}$} \& |[lcs header]|{B} \& |[lcs header]|{D} \& |[lcs header]|{C} \& |[lcs header]|{A} \& |[lcs header]|{B} \& |[lcs header]|{A}\\
  |[lcs header]|{$\boldsymbol{\epsilon}$} \& 0\&0\&0\&0\&0\&0\&0\\
  |[lcs header]|{A} \&0\&0\&0\&0\&1\&1\&1\\
  |[lcs header]|{B} \&0\&1\&1\&1\&1\&2\&2\\
  |[lcs header]|{C} \&0\&1\&1\&2\&2\&2\&2\\
  |[lcs header]|{B} \&0\&1\&1\&2\&2\&3\&3\\
  |[lcs header]|{D} \&0\&1\&2\&2\&2\&3\&3\\
  |[lcs header]|{A} \&0\&1\&2\&2\&3\&3\&4\\
  |[lcs header]|{B} \&0\&1\&2\&2\&3\&4\&4\\
};
\draw[trace move] (lcs-9-8.north) -- (lcs-8-8.south);
\draw[trace move] (lcs-8-8.north west) -- (lcs-7-7.south east);
\draw[trace move] (lcs-7-7.north) -- (lcs-6-7.south);
\draw[trace move] (lcs-6-7.north west) -- (lcs-5-6.south east);
\draw[trace move] (lcs-5-6.west) -- (lcs-5-5.east);
\draw[trace move] (lcs-5-5.north west) -- (lcs-4-4.south east);
\draw[trace move] (lcs-4-4.west) -- (lcs-4-3.east);
\draw[trace move] (lcs-4-3.north west) -- (lcs-3-2.south east);
\foreach \r/\c/\ch in {8/8/A,6/7/B,5/5/C,4/3/B}{
  \node[chosen,fit=(lcs-\r-\c)] {};
}
\node[font=\scriptsize,text=AlgoGray,align=left,right=4mm of lcs-7-8]{
$(7,6)\uparrow(6,6)\nwarrow_{\!A}(5,5)$\\
$\uparrow(4,5)\nwarrow_{\!B}(3,4)\leftarrow(3,3)$\\
$\nwarrow_{\!C}(2,2)\leftarrow(2,1)\nwarrow_{\!B}(1,0)$};
\end{tikzpicture}
\vspace{1mm}
\[
\text{selected backward: }\texttt{A,B,C,B}
\qquad\xrightarrow{\text{reverse}}\qquad
\textcolor{green!45!black}{\textbf{\texttt{BCBA}}}
\]
\end{frame}
\begin{frame}{복원 결과 검증}
\begin{itemize}\item \texttt{BCBA}는 $X=\texttt{ABCBDAB}$의 subsequence\item \texttt{BCBA}는 $Y=\texttt{BDCABA}$의 subsequence\item 길이 4는 $dp[7][6]$과 일치\end{itemize}
\end{frame}
\begin{frame}{LCS Space Optimization}
\optional[선택 심화]\hfill
\begin{itemize}\item 길이만 필요: previous/current 두 row → $\Theta(\min(m,n))$ auxiliary space\item time은 여전히 $\Theta(mn)$\item full table을 버리면 단순 backtracking 정보도 사라짐\end{itemize}
\begin{block}{Appendix preview}Hirschberg는 선형 공간으로 sequence를 복원하지만 추가 divide-and-conquer가 필요하다.\end{block}
\vfill\muted{LCS는 2차원 prefix state였다. 다음 문제에서는 1차원 suffix state와 전체 답을 분리한다.}
\end{frame}
